\section{习题课 \#10（11月15日）}

\subsection{矩阵的等价与线性映射的表示}

设 $T:V\to W$ 是线性映射。取 $V$ 的基
\begin{equation}
 \mathcal A=(\alpha_1,\ldots,\alpha_n),
 \qquad
 \widetilde{\mathcal A}=(\widetilde\alpha_1,\ldots,\widetilde\alpha_n),
\end{equation}
以及 $W$ 的基
\begin{equation}
 \mathcal B=(\beta_1,\ldots,\beta_m),
 \qquad
 \widetilde{\mathcal B}=(\widetilde\beta_1,\ldots,\widetilde\beta_m).
\end{equation}
若
\begin{equation}
 [v]_{\mathcal A}=Q[v]_{\widetilde{\mathcal A}},
 \qquad
 [w]_{\mathcal B}=P[w]_{\widetilde{\mathcal B}},
\end{equation}
且 $T$ 在两组基下的矩阵分别为 $A$、$\widetilde A$，则
\begin{equation}
 \widetilde A=P^{-1}AQ.
\end{equation}
因此，同一线性映射在不同定义域和值域基下的矩阵彼此等价。

\begin{definition}
同型矩阵 $A,B\in\F^{m\times n}$ 称为等价，如果存在可逆矩阵
\begin{equation}
 P\in\GL_m(\F),
 \qquad Q\in\GL_n(\F)
\end{equation}
使
\begin{equation}
 B=PAQ.
\end{equation}
\end{definition}

\begin{theorem}
两个同型矩阵等价当且仅当它们的秩相同。每个秩为 $r$ 的 $m\times n$ 矩阵都等价于
\begin{equation}
 \begin{pmatrix}I_r&0\\0&0\end{pmatrix}.
\end{equation}
\end{theorem}

\subsection{相似与线性变换}

当 $V=W$ 且定义域和值域同时更换为同一组基时，左右变换互为逆矩阵。

\begin{definition}
方阵 $A,B\in\F^{n\times n}$ 称为相似，如果存在 $Q\in\GL_n(\F)$ 使
\begin{equation}
 B=Q^{-1}AQ.
\end{equation}
相似矩阵表示同一个线性变换在不同基下的矩阵。
\end{definition}

\begin{proposition}
对任意 $Q\in\GL(V)$，映射
\begin{equation}
 \Phi_Q:\End(V)\to\End(V),
 \qquad
 \Phi_Q(T)=Q^{-1}TQ
\end{equation}
是代数自同构。也就是说，
\begin{align}
 \Phi_Q(aS+bT)&=a\Phi_Q(S)+b\Phi_Q(T),\\
 \Phi_Q(ST)&=\Phi_Q(S)\Phi_Q(T),\\
 \Phi_Q(I)&=I.
\end{align}
其逆为 $\Phi_{Q^{-1}}$。
\end{proposition}

\begin{remark}
矩阵等价只保留秩，而矩阵相似保留线性变换的迭代结构，因而分类精细得多。相同特征多项式并不足以保证相似，例如
\begin{equation}
 \begin{pmatrix}0&1\\0&0\end{pmatrix}
 \quad\text{与}\quad
 \begin{pmatrix}0&0\\0&0\end{pmatrix}
\end{equation}
特征多项式同为 $\lambda^2$，但秩不同，所以不相似。
\end{remark}

\subsection{相似不变量}

若 $B=Q^{-1}AQ$，则
\begin{align}
 \rank B&=\rank A,\\
 \det B&=\det A,\\
 \tr B&=\tr A,\\
 \det(\lambda I-B)&=\det(\lambda I-A),\\
 f(B)&=Q^{-1}f(A)Q
\end{align}
对任意多项式 $f$ 成立。特别地，相似关系保持以下性质：
\begin{enumerate}
  \item 可逆性；
  \item 幂等性 $A^2=A$；
  \item 幂零性 $A^k=0$；
  \item 最小多项式与特征多项式；
  \item 各次幂的秩与核维数。
\end{enumerate}

\begin{remark}
相似一般不保持对称性、Hermite 性或正交性，因为这些性质还依赖于所选内积和基是否标准正交。若相似变换矩阵是正交或酉矩阵，则这些带内积的结构才会保持。
\end{remark}

\subsection{特征值与特征空间}

对 $A\in\F^{n\times n}$，若存在非零向量 $v$ 满足
\begin{equation}
 Av=\lambda v,
\end{equation}
则 $\lambda$ 称为特征值，$v$ 称为相应特征向量。特征空间为
\begin{equation}
 V_\lambda=\Ker(\lambda I-A).
\end{equation}
特征值由
\begin{equation}
 \chi_A(\lambda)=\det(\lambda I-A)=0
\end{equation}
确定。

\begin{theorem}[不同特征空间的直和性]
设 $\lambda_1,\ldots,\lambda_s$ 两两不同。若
\begin{equation}
 v_i\in V_{\lambda_i},
 \qquad i=1,\ldots,s,
\end{equation}
且每个 $v_i\neq0$，则 $v_1,\ldots,v_s$ 线性无关。更一般地，
\begin{equation}
 V_{\lambda_1}+\cdots+V_{\lambda_s}
\end{equation}
是直和。
\end{theorem}

\begin{proof}
设
\begin{equation}
 v_1+\cdots+v_s=0.
\end{equation}
对等式作用
\begin{equation}
 \prod_{j\neq i}(A-\lambda_jI),
\end{equation}
除第 $i$ 项外其余项都消失，得到
\begin{equation}
 \prod_{j\neq i}(\lambda_i-\lambda_j)v_i=0.
\end{equation}
因特征值两两不同，系数非零，故 $v_i=0$。
\end{proof}

\begin{corollary}
若 $A$ 在 $\F$ 中有 $n$ 个互不相同的特征值，则 $A$ 可对角化。
\end{corollary}

\begin{definition}
特征值 $\lambda$ 在特征多项式中的重数称为代数重数，记为 $m_a(\lambda)$；特征空间维数
\begin{equation}
 m_g(\lambda)=\dim\Ker(\lambda I-A)
\end{equation}
称为几何重数。
\end{definition}

\begin{theorem}
总有
\begin{equation}
 1\leq m_g(\lambda)\leq m_a(\lambda).
\end{equation}
矩阵 $A$ 可对角化当且仅当特征多项式在 $\F$ 中完全分裂，且每个特征值的几何重数等于代数重数。
\end{theorem}

\subsection{两个特征值计算例}

\begin{example}[可对角化矩阵及其幂]
原稿中的第一个矩阵为
\begin{equation}
 A_1=\begin{pmatrix}
 -1&3&-1\\
 -3&5&-1\\
 -3&3&1
 \end{pmatrix}.
\end{equation}
直接计算得
\begin{equation}
 \chi_{A_1}(\lambda)
 =\det(\lambda I-A_1)
 =(\lambda-2)^2(\lambda-1).
\end{equation}
当 $\lambda=2$ 时，可取两个线性无关特征向量
\begin{equation}
 v_1=(1,1,0)^T,
 \qquad
 v_2=(1,0,-3)^T.
\end{equation}
当 $\lambda=1$ 时，可取
\begin{equation}
 v_3=(1,1,1)^T.
\end{equation}
因此
\begin{equation}
 P=(v_1,v_2,v_3),
 \qquad
 A_1=P\diag(2,2,1)P^{-1},
\end{equation}
并且
\begin{equation}
 A_1^m=P\diag(2^m,2^m,1)P^{-1}.
\end{equation}
这比逐次相乘更适合计算高次幂。
\end{example}

\begin{example}[实矩阵的复特征值与实分块]
令
\begin{equation}
 A_2=\begin{pmatrix}
 4&-5&7\\
 1&-4&9\\
 -4&0&5
 \end{pmatrix}.
\end{equation}
计算得
\begin{equation}
 \chi_{A_2}(\lambda)
 =(\lambda-1)\bigl((\lambda-2)^2+3^2\bigr).
\end{equation}
所以在 $\R$ 上只有特征值 $1$，可取特征向量
\begin{equation}
 v_1=(1,2,1)^T,
\end{equation}
矩阵不能在 $\R$ 上对角化。在 $\C$ 上另有
\begin{equation}
 \lambda_2=2+3\ii,
 \qquad
 \lambda_3=2-3\ii,
\end{equation}
可取
\begin{equation}
 v_2=(3-3\ii,5-3\ii,4)^T,
 \qquad
 v_3=\overline{v_2}.
\end{equation}
于是 $A_2$ 在 $\C$ 上可对角化。

也可在实数域内令
\begin{equation}
 u=\Re v_2=(3,5,4)^T,
 \qquad
 w=\Im v_2=(-3,-3,0)^T.
\end{equation}
则 $\Span\{u,w\}$ 是实不变子空间，并且在基 $(u,w)$ 下限制矩阵为
\begin{equation}
 \begin{pmatrix}2&3\\-3&2\end{pmatrix}
\end{equation}
或按基次序不同得到其转置型。故 $A_2^m$ 可由一个一维块 $1^m$ 和该二维旋转--伸缩块的 $m$ 次幂计算。
\end{example}

\subsection{Gershgorin 圆盘定理}

\begin{theorem}[Gershgorin]
设 $A=(a_{ij})\in\C^{n\times n}$，并定义圆盘
\begin{equation}
 D_i=\set{z\in\C:
 |z-a_{ii}|\leq\sum_{j\neq i}|a_{ij}|}.
\end{equation}
则每个特征值都属于至少一个圆盘：
\begin{equation}
 \sigma(A)\subseteq\bigcup_{i=1}^nD_i.
\end{equation}
\end{theorem}

\begin{proof}
设 $Ax=\lambda x$，$x\neq0$，取 $k$ 使
\begin{equation}
 |x_k|=\max_i|x_i|.
\end{equation}
第 $k$ 行给出
\begin{equation}
 (\lambda-a_{kk})x_k
 =\sum_{j\neq k}a_{kj}x_j.
\end{equation}
因此
\begin{equation}
 |\lambda-a_{kk}|
 \leq\sum_{j\neq k}|a_{kj}|
 \frac{|x_j|}{|x_k|}
 \leq\sum_{j\neq k}|a_{kj}|,
\end{equation}
即 $\lambda\in D_k$。
\end{proof}

\begin{remark}
若若干 Gershgorin 圆盘与其余圆盘不相交，则该连通部分中包含的特征值个数等于其中圆盘的个数，按代数重数计算。这一加强结论可由连续变形到对角矩阵得到。
\end{remark}

\subsection{上三角化与 Schur 分解}

\begin{theorem}
若 $A\in\C^{n\times n}$，则存在可逆矩阵 $P$ 使
\begin{equation}
 P^{-1}AP
\end{equation}
为上三角矩阵。其对角元正是 $A$ 的特征值，按代数重数排列。
\end{theorem}

\begin{proof}
复数域上特征多项式有根。取一个特征向量作为新基的第一个向量，则矩阵具有分块形式
\begin{equation}
 \begin{pmatrix}\lambda&*\\0&A_1\end{pmatrix}.
\end{equation}
对 $A_1$ 归纳即可。
\end{proof}

\begin{theorem}[Schur 酉三角化]
对任意 $A\in\C^{n\times n}$，存在酉矩阵 $U$ 使
\begin{equation}
 U^*AU=T
\end{equation}
为上三角矩阵。
\end{theorem}

\begin{remark}
实矩阵未必有实特征值，因此未必能在 $\R$ 上三角化。但存在正交矩阵 $Q$，使 $Q^TAQ$ 为实拟上三角矩阵，其对角块为 $1\times1$ 实特征值块或对应一对共轭复特征值的 $2\times2$ 实块。这是实 Schur 分解。
\end{remark}
